\renewcommand{\ElicitationDate}{2026.10.22}

\section{Elicitation \ElicitationDate}

\begin{tabular}{ll}
\textbf{Date:} & \ElicitationDate \\
\textbf{Consultant:} & Hector Hill / Gy'ax \\
\textbf{Elicitors:} & Zahra, Andrew, Matthias \\
\textbf{Location:} & TFS elicitation room \\
\textbf{Context:} & Fifth elicitation for LING 431/531 Field Methods \\
\textbf{Audio:} & \texttt{October\_22\_2026\_HH.WAV} \\
\end{tabular}

\subsection{Zahra [2:30-35:40]}

[2:40]
\begin{exe}
    \ex \glll
    K’i’yhl t’imḵ’aax g̱a’negwihl ha’niitx̱ooḵxw ii k’i’yhl sdo’ohl t’imḵ’aax g̱ale’yhlit. \\
    k’i’y=hl t’imḵ’aax g̱a-’nekw=hl *(ha’niitx̱ooḵxw) ii *k’i’y=hl sdo’o=hl t’imḵ’aax g̱a-le’y=hl-it. \\
    one=\Cn{} surface-armspan \Dist-long=\Cn{} \Ins-on-eat-\Val{} one half=\Cn{}surface-armspan \Dist-wide=\Cn{} \Ins-on-eat-\Val=\Cn-3.\Ii \\
    \glt `This table is an armspan long and half an armspan wide.' \hfill HH.\ElicitationDate{} (SF)
\end{exe}

Should double check whether \textit{-t} is \Sx{} or 3.\Ii; Henry thought \Sx, might be because of the \Ccnj.

[04:30]
\begin{exe}
    \ex[]{\glll
    K’i’yhl t’imḵ’aax g̱a’negwihl g̱an k’i’yhl sdo’ohl t’imḵ’aax g̱ale’yhl ha’niitx̱ooḵxw. \\
    k’i’y=hl t’im-ḵ’aax g̱a-’nakw=hl g̱an k’i’y=hl sdo’o=hl t’im-ḵ’aax g̱ale’yhl ha’niitx̱ooḵxw. \\
    one=\Cn{} surface-armspan \Dist-long=\Cn{} \Pcnj{} one half=\Cn{} surface-armspan \Dist-wide=\Cn{} table \\
    \glt \textit{Intended}: "The table is an armspan long and half an armspan wide." \hfill HH.\ElicitationDate{} (SF) \\
    Comment: "Without the \textit{ii} you're not talking about the same table ... when you put \textit{g̱an} that's added on to another table."}
\end{exe}

\begin{exe}
    \ex[]{\glll
    K’i’yhl t’imḵ’aax g̱a’negwihl ha’niitx̱ooḵxw g̱an k’i’yhl sdo’ohl t’imḵ’aax g̱ale’yhlit. \\
    k’i’y=hl t’im-ḵ’aax g̱a-’nakw=hl g̱an k’i’y=hl sdo’o=hl t’im-ḵ’aax g̱ale’yhl ha’niitx̱ooḵxw. \\
    one=\Cn{} surface-armspan \Dist-long=\Cn{} table \Pcnj{} one half=\Cn{} surface-armspan \Dist-wide=\Cn{} table \\
    \glt "The table is an armspan long and half an armspan wide." \hfill HH.\ElicitationDate{} (SF) 
    } 
\end{exe}

The judgment on this last one was a bit inconclusive.

[9:53]
\begin{exe}
    \ex \glll
    K’i’y=hl t’imḵ’aax g̱an k’i’y=hl sdo’ohl (*t’imḵ’aax) g̱a-’nekw=hl ha’niitx̱ooḵxw. \\
    k’i’y=hl t’im-ḵ’aax g̱an k’i’y=hl sdo’o=hl t’im-ḵ’aax g̱a-’nakw=hl ha’niitx̱ooḵxw \\
    one=\Cn{} surface-armspan \Pcnj{} one=\Cn{} half=\Cn{} surface-armspan \Dist-long=\Cn{} table \\
    \glt "The table is an armspan and a half wide." \hfill HH.\ElicitationDate{} (SF) \\
    Comment: Having \textit{t’imḵ’aax} in the second conjunct is repetitive. \\
    Comment: We're talking about the same table.
\end{exe}

[13:30]
\begin{exe}
    \ex \glll
    K’i’yhl t’imḵ’aax g̱a’nakwihl ha’niitx̱ooḵxw ii k’i’yhl sdo’ohl t’imḵ’aax ha’niitx̱ooḵxw. \\
    k’i’y=hl t’im-ḵ’aax g̱a-’nekw=hl ha’niitx̱ooḵxw ii k’i’y=hl sdo’o=hl t’im-ḵ’aax ha’niitx̱ooḵxw. \\
    one=\Cn{} surface-armspan \Pcnj{} one=\Cn{} half=\Cn{} surface-armspan \Dist-long=\Cn{} table \\
    \glt "The table is an armspan long and the [other?] table is half an armspan." \hfill HH.\ElicitationDate{} (VF) \\
    Comment: Mentioning two \textit{ha’niitx̱ooḵxw}. If you replace the second one with -\textit{it}, you're talking about the same table.
\end{exe}

[14:22]
\begin{exe}
    \ex \glll
    Xwsdins se-se’e=hl g̱a-’nakw=hl ha’niitx̱ooḵxw ‘ii gwile’l si-se’   g̱a-le’y=hl-it … ‘ii  \\ 
    \\ 
    Five 	 PL-foot=CN g̱a-long=CN (table)	    and three PL-foot=CN g̱a-wide=CN (table) \\
    \glt "This table is five feet long and three feet wide. … therefore [fifteen] feet …." \hfill HH.\ElicitationDate{} (VF)
\end{exe}

[15:39]
\begin{exe}
    \ex \glll
    ... sag̱ayt tunhl g̱asg̱oohl ha’niitx̱ooḵxw \\
    {} sag̱ayt t=xʷin=hl g̱asg̱oohl=\Cn{} table \\
    {} together \Pn=\Dem.\Prox=\Cn{} \Mass=\Cn{} table \\
    \glt "... is [together, in all] the size of the table." \hfill HH.\ElicitationDate{} (VF)
\end{exe}

[18:37]
\begin{exe}
    \ex \glll
    gwile'l g̱abiihl ha'niit'aa dim 'nii t'aat lax ha’niitx̱ooḵxw \\
    gwile'l g̱abii=hl ha-'nii-t'aa dim 'nii t'aa-t lax ha’-nii-tx̱ooḵ-xw \\
    two \Cnt=\Cn{} \Ins-on-sit \Prosp{} on sit-3.\Sx{} on \Ins-on-eat-\Val 
    \\
    \glt "two chairs that they'll sit on at the table" (orig. tr.) \hfill HH.\ElicitationDate{} (VF)
\end{exe}

Would be good to double check this translation; may also be a fragment. Also check which -\textit{t} is on \textit{t'aat}.

[20:40]
\begin{exe}
    \ex \glll Gwile'lhl sise’ehl g̱a’nagwihl ha’niitx̱ooḵxw ii gilbilhl sise’e g̱ale’yhlit, ii sag̱ayt ḵ’oolt g̱asg̱oot. \\
    gwile'l=hl si-se’e=hl g̱a-’nakw=hl ha’niitx̱ooḵxw ii gilbil=hl si-se’e=hl   g̱a-le’y=hl-it, ii sag̱ayt ḵ'oolt g̱asg̱oo-t. \\
    three=\Cn{} \Pl-foot=\Cn{} \Distr-long=\Cn{} table and two \Pl-foot=\Cn{} \Distr-long=\Cn-3.\Ii, \Ccnj{} together six \Mass-3.\Ii \\ 
    \glt `The table is three feet long and two feet wide; together six [feet in area].'\hfill HH.\ElicitationDate{} (SF)
\end{exe}

Might need to know more about the syntax of such expressions to know if the last -\textit{t} is 3.\I{} or 3.\Ii. The suffix itself is also not clearly audible so would be good to double check that too. \\

There are a couple points when HH seems to be adding rather than multiplying with \textit{sag̱ayt} `together,' but other times when he seems to accept multiplication. 

[21:59]
\begin{exe}
    \ex \glll
    sag̱ayt ’nithl g̱asg̱oot \\
    sag̱ayt ’nit=hll g̱asg̱oo-t \\
    together 3.\Iii=\Cn{} \Mass-3.\Ii \\
    \glt "it's added on; total" \hfill HH.\ElicitationDate{} (VF)
\end{exe}

[23:30]
\begin{exe}
    \ex \glll g̱ay k’aa ‘wii ’nekw=hl ha’niitx̱ooḵxw ehl (\textbf{win} \textbf{sdo’ohl}) g̱ale’yhlit. \\
    g̱ay k’aa ‘wii ’nekw=hl ha’niitx̱ooḵxw a[-t]=hl (\textbf{win} \textbf{sdo’o=hl}) g̱a-le’y=hl-t \\
    \Cntr{} exceedingly big long=\Cn{} table \Prep{}[-3.\Ii]=\Cn{} \Comp{} half=\Cn{} \Distr-long=\Cn{}-3.\Ii \\ 
    \glt `The table is longer than it is wide (\textbf{by half}).' \hfill HH.\ElicitationDate{} (VF)
\end{exe}

HH seems to say \textit{sdo'okshl} for \textit{sdo'ohl} `half' at times but that seems to be competition from the separate lexical item \textit{sdo'oks} `beside.' Not sure if they're related of course. He corrects himself to \textit{sdo'ohl}.

[27:37]
\begin{exe}
    \ex \glll
    g̱ay k’aa ‘wii g̱a’negwihl ha’niitx̱ooḵxw ahl g̱ale’yhlit, ii ḵ’ap nee dii ‘wii (*g̱a)‘nakwt. \\
    g̱ay k’aa ‘wii g̱a-’nekw=hl ha’niitx̱ooḵxw a[-t]=hl g̱a-le’y=hl-it, ii ḵ’ap nee dii ‘wii g̱a-‘nakw-it. \\
    \Cntr{} more big \Distr-long=\Cn{} table \Prep{}[-3.\Ii]=\Cn{} \Distr-long=\Cn{}-3.\Ii, and \Verum{} \Neg{} \Foc{} big \Distr-long-3.\Ii \\
    \glt `The table is longer than it is wide, but actually it is not long.'
    \hfill HH.\ElicitationDate{} (SF)
\end{exe}

Wasn't totally clear if the \textit{g̱a}- in the first \textit{g̱a’negwihl} should be there. I think there was also something interesting here about \textit{g̱a’negwihl} \Distr-long=\Cn{}  vs. \textit{g̱a’nakwt} \Distr-long-3.\Ii{} vs. \textit{g̱a’negwit} \Distr-long-\Sx{}, with implications about whether you're talking about the same or another table. But will want to clarify this a bit. 

\begin{exe}
    \ex \glll
    g̱ay k’aa ‘wii ’negwit \\
    g̱ay k’aa ‘wii ’nakw-it \\
    \Cntr{} more big long-3.\Sx? \\
    \glt "it (this table) is longer" \hfill HH.\ElicitationDate{} (VF) \\
    Comment: "There's two different tables sitting there, and you're pointing to one."
\end{exe}


\subsection{Andrew []}